\chapter*{\lblHowToUse}
\phantomsection
\addcontentsline{toc}{chapter}{\lblHowToUse}
\markboth{\lblHowToUse}{\lblHowToUse}

\noindent
Ta książka jest złożona pod ekran, nie pod półkę. Strona ma format A4,
dokument jest jednostronny, a każda decyzja projektowa poniżej zakłada ten sam
obraz: \textbf{PDF otwarty na jednej połowie ekranu, edytor otwarty na
drugiej}, a gdzieś pod edytorem terminal. Jeśli czytasz ją w pociągu bez
laptopa, dostaniesz argument i ominiesz sedno, bo sednem jest to, co dzieje
się w edytorze.

\section*{Pętla}

Każdy rozdział czyta się tak samo, a pętla jest na tyle krótka, że po trzech
rozdziałach staje się nawykiem:

\begin{enumerate}[leftmargin=1.6em, itemsep=3pt]
  \item Przeczytaj sekcję. Kończy się listingiem, a listing jest
        \emph{plikiem}: jego ścieżka jest wydrukowana nad nim, względem
        katalogu głównego repozytorium, i to ten plik otwierasz w edytorze.
  \item Uruchom plik. Każdy listing w tej książce działa tak, jak go
        wydrukowano, na wersjach ze strony tytułowej, a build repozytorium
        uruchamia go ponownie przy każdej zmianie. Jeśli u ciebie nie działa,
        jedno z was ma złą wersję, a strona tytułowa mówi które.
  \item Zrób ćwiczenie. Ćwiczenie to plik startowy i test obok niego. Spraw,
        żeby test przeszedł, a potem przeczytaj, co test sprawdzał i dlaczego.
  \item Idź dalej. Następna sekcja zakłada, że zrobiłeś poprzednie trzy kroki.
\end{enumerate}

\noindent
\mermaidfig{reading-loop}{Pętla: przeczytaj sekcję w PDF-ie, otwórz i uruchom jej listing w edytorze, spraw, żeby test ćwiczenia przeszedł, idź dalej}{pętla czytania}

Oto cała pętla, raz, zanim będzie się czego z niej uczyć. Listing poniżej
jest plikiem; otwórz go.

\pyfile{code/ch00/loop.py}{Pierwszy listing i kształt każdego następnego}{lst:ch00-loop}

\noindent
Uruchomiony z katalogu głównego repozytorium drukuje to, co poniżej \dash{}
również wygenerowane, a nie przepisane: ta linia jest tym, co przypięty
interpreter wydrukował przy ostatnim buildzie książki:

\transcript{ch00-loop}

\noindent
A oto ćwiczenie, które do niego należy. Ma numer 0.1, bo ten rozdział nie
jest rozdziałem; pierwsze prawdziwe to Ćwiczenie~1.1.

\exercise{e00_01_hello}{Spraw, żeby pierwszy test przeszedł}{%
Otwórz plik startowy. Zawiera jedną funkcję, która rzuca
\code{NotImplementedError}, a test obok niej nie przechodzi, dopóki funkcja
nie zwróci tego, czego test oczekuje. Spraw, żeby przeszedł. Potem uruchom
test ponownie z \code{-v} i przeczytaj, co wydrukował pytest, bo sposób, w
jaki pytest raportuje niepowodzenie, to pierwsza rzecz w tej książce, która
nie ma odpowiednika w .NET wartego nazwania.}

\noindent
Każde ćwiczenie w książce drukuje te trzy linie: plik, który otwierasz, plik,
który decyduje, czy skończyłeś, i jedno polecenie, które go uruchamia. Plik
startowy każdego ćwiczenia celowo nie przechodzi swojego testu, a build
repozytorium sprawdza, że tak jest \dash{} ćwiczenie, którego plik startowy
już przechodzi, nie jest ćwiczeniem.

\section*{Czym jest rozdział}

Czternaście rozdziałów, a każdy trzyma się tego samego kontraktu:

\begin{itemize}[leftmargin=1.4em, itemsep=2pt]
  \item \textbf{Jeden argument}, postawiony w pierwszym akapicie i spłacony w
        ostatnim. Rozdział jest tezą o tym, gdzie nawyk z C\# przestaje
        działać i co go zastępuje, a nie wycieczką po bibliotece.
  \item \textbf{Poniżej trzech tysięcy słów} prozy, listingów nie licząc.
        Build je liczy i odrzuca rozdział, który ma za dużo. Trzy tysiące
        słów to jedno posiedzenie; książka to czternaście posiedzeń.
  \item \textbf{Tabela mapowania}, C\# po lewej i Python po prawej, bo
        czytelnik ma już pojęcie i potrzebuje nazwy.
  \item \textbf{Listingi, które są plikami}, każdy uruchomiony na przypiętych
        wersjach, zanim został wydrukowany. Tam, gdzie listing nie został
        uruchomiony, mówi o tym w ramce, a build liczy ramki.
  \item \textbf{Pułapki}, wywoływane, zanim zostaną nazwane. Tam, gdzie nawyk
        z C\# daje Pythona, który działa i jest zły, rozdział pozwala ci go
        najpierw napisać. Dodatek~\ref{app:traps} jest katalogiem.
  \item \textbf{Ćwiczenia} z plikiem startowym i testem, jak wyżej.
  \item \textbf{Niezależność.} Rozdział nazywa to, co zakłada z wcześniejszych,
        i daje się czytać osobno; każdy napisano tak, żeby dało się go
        opublikować osobno.
\end{itemize}

\section*{Ramki}

\begin{center}
\small
\begin{tabularx}{\linewidth}{@{}lX@{}}
\toprule
\textbf{\lblNote} & Coś wartego wiedzy, czego argument nie potrzebuje. \\
\textbf{\lblWarning} & Coś, co będzie cię kosztować popołudnie. \\
\textbf{\lblCsharp} & Tłumaczenie. Strona C\# zdania, które właśnie
  przeczytałeś, i nic więcej; czytelnik bez .NET za sobą może pominąć każdą z
  tych ramek i nie stracić treści. \\
\textbf{\lblTrap} & Błędne przekonanie, nazwane po tym, jak w nie wszedłeś.
  Każde jest w Dodatku~\ref{app:traps}. \\
\textbf{\lblVersion} & Wersja zapoznawcza, eksperymentalna albo taka, która
  może się zmienić. Rzadkie w książce o języku; częste w jej dwóch ostatnich
  rozdziałach. \\
\textbf{\lblVerify} & Listing, którego nie uruchomiono na przypiętych
  wersjach. Nic nie trafia do czytelnika z taką ramką. \\
\textbf{\lblExercise} & Plik startowy i test. \\
\textbf{\lblProject} & Etap \code{trace-assert}, który buduje ta sekcja.
  Zobacz Wprowadzenie. \\
\bottomrule
\end{tabularx}
\end{center}

\section*{Dwie zasady, których ta książka się trzyma}

\begin{note}
\textbf{Każdy listing został uruchomiony albo mówi, że nie został.} Kod w
katalogu \code{code/} jest osobnym projektem z plikiem blokady, a build
repozytorium instaluje dokładnie wersje ze strony tytułowej i uruchamia każdy
listing i każdy test ćwiczenia przy każdej zmianie. Listing, który przez to
nie przeszedł, nosi ramkę \emph{\lblVerify}, a build drukuje, ile ich jest.
Kiedy ta liczba wynosi zero, książka ma prawo do swojej strony tytułowej.
\end{note}

\begin{note}
\textbf{To, czego nie zmierzono, jest oznaczone jako osąd.} Tam, gdzie
książka mówi, że jedno jest szybsze od drugiego, jest skrypt w katalogu
\code{code/measure/}, który wyprodukował liczbę, a liczba trafia na stronę
mechanicznie, a nie przez przepisanie. Tam, gdzie żaden skrypt jeszcze nie
został uruchomiony, zdanie to mówi, a tabela, która miałaby zawierać liczby,
zostaje pusta. Puste tabele nie są przeoczeniem; są mechanizmem.
\end{note}

\section*{Trzy wydania}

Książka ukazuje się w trzech częściach, a każda jest kompletna sama w sobie.
Wersja~0.1 to części~I i~II \dash{} środowisko uruchomieniowe, narzędzia i
język. Wersja~0.2 dodaje części~III i~IV \dash{} współbieżność i wdrażanie.
Wersja~1.0 dodaje część~V i dodatki. Rozdział, którego jeszcze nie
napisano, drukuje czerwoną ramkę, która o tym mówi, a build je również
liczy.

\section*{Jeśli już piszesz w Pythonie}

Przeczytaj najpierw Dodatek~\ref{app:traps} i bądź szczery co do wpisów,
których nie znałeś. Potem i tak przeczytaj Rozdziały~\ref{ch:asyncio}
i~\ref{ch:testing}, bo to tam zaplecze C\# produkuje Pythona, który działa i
jest zły, a rok pisania w Pythonie nie zawsze wystarcza, żeby to zauważyć.

\section*{Dwa języki, jedno źródło}

Ta książka istnieje po angielsku i po polsku, budowana z jednego
repozytorium. Oba wydania dzielą rozdziały, listingi, ćwiczenia i numerację;
różni się tylko proza. Sprawdzenie przy każdym buildzie dowodzi, że oba mają
tę samą strukturę i te same liczby, więc poprawka naniesiona w jednym wydaniu
nie może po cichu ominąć drugiego. Angielski jest pierwszym językiem książki:
kod, identyfikatory i każdy komentarz w listingu są po angielsku w obu
wydaniach, bo czytelnik jest uczony czytania angielskiego kodu. Tam,
gdzie polski termin nie ma ustalonej formy w realnym użyciu, wydanie
polskie zachowuje angielskie słowo i odmienia je, i mówi o tym w tym
miejscu.
